\documentclass[12pt,a4paper]{article}
\usepackage[margin=2.5cm]{geometry}
\usepackage{amsmath}
\usepackage{hyperref}
\usepackage{parskip}

\title{\textbf{Code Documentation}\\[0.3em]
\large Part V(g) --- Modelling Report\\
Asymmetric Sentiment Effects in Earnings Calls}
\author{NLP Course --- Group Project}
\date{}

\begin{document}
\maketitle

\section{Package Structure}

The codebase is organized as an installable Python package (\texttt{nasdaq-nlp})
under \texttt{src/nasdaq\_nlp/}, following the src-layout convention.
This prevents accidental imports of the uninstalled source directory.

\begin{verbatim}
src/nasdaq_nlp/
├── config.py           All output file paths as Path constants
├── data/
│   ├── loader/         Transcript scanning (both dataset sources)
│   ├── metadata/       Header parsing, event trading day assignment
│   └── market/         Price download and return computation
├── preprocessing/      Text cleaning, section split, tokenization
├── features/
│   ├── lexicon.py      LM sentiment rates (full, pres, Q&A)
│   ├── tfidf.py        TF-IDF fitting and transform
│   └── embeddings.py   Word2Vec training and document pooling
├── models/
│   ├── regression.py   OLS, Ridge, RF, MLP, Wald test, time-series CV
│   └── classifiers.py  NB, LogReg, Tree, RF, MLP classifiers
├── evaluation/         OOS R², MAE, AUC utilities
└── visualization.py    All benchmark plots
\end{verbatim}

\section{Module Documentation}

Every module begins with a docstring explaining its purpose and public API.
Every public function includes \texttt{Parameters} and \texttt{Returns} sections.
Example from \texttt{data/loader/main.py}:

\begin{verbatim}
def scan_transcripts(
    dataset_dir: Path = DATASET_DIR,
    ect_dir: Path = ECT_DATASET_DIR,
    original: bool = True,
    ect: bool = True,
) -> list[TranscriptRecord]:
    """Scan both transcript datasets and return a combined list.

    Parameters
    ----------
    dataset_dir : Path
        Root of the Thomson Reuters dataset.
    ect_dir : Path
        Root of the Kaggle cleaned_ECTs dataset.
    original : bool
        Include original-source transcripts (default True).
    ect : bool
        Include ECT-source transcripts (default True).

    Returns
    -------
    list[TranscriptRecord]
        Sorted by (ticker, year, quarter).
    """
\end{verbatim}

\section{Setup and Installation}

\begin{verbatim}
# Install UV package manager (if not already installed)
curl -LsSf https://astral.sh/uv/install.sh | sh

# Install all dependencies and the nasdaq_nlp package
uv sync
\end{verbatim}

\section{Running the Pipeline}

\begin{verbatim}
make pipeline     # full pipeline: metadata -> returns -> features -> models
make serve        # launch Streamlit dashboard
make lint         # ruff linter check
make format       # auto-fix formatting with ruff
make clean        # delete outputs/ (keeps dataset/ intact)
\end{verbatim}

Individual pipeline steps can be re-run in isolation.
Each step reads its input CSV, applies a single transformation, and writes its
output CSV to \texttt{outputs/processed/}.

\section{Reproducibility}

\begin{itemize}
  \item Dependency versions are pinned in \texttt{uv.lock}.
  \item All sklearn models use \texttt{random\_state=42}.
  \item Word2Vec is the only non-deterministic step; its output is cached after
        the first run so subsequent runs are deterministic.
  \item The full analysis can be reproduced from scratch via:
\begin{verbatim}
uv sync
make pipeline
jupyter nbconvert --execute notebooks/03_modeling.ipynb
jupyter nbconvert --execute notebooks/04_results.ipynb
\end{verbatim}
\end{itemize}

\section{Code Quality}

Static analysis and formatting are enforced by \texttt{ruff}:
\begin{verbatim}
make lint     # ruff check + format --check on src/
make format   # ruff check --fix + format src/
\end{verbatim}

All modules in \texttt{src/nasdaq\_nlp/} pass ruff without warnings.

\section{Repository}

Full source code, notebooks, and documentation:
\url{https://github.com/Tee-Time-Software-Solutions/NASDAQ-NLP}

\end{document}
